% ---- pg-036 （书页 24） ----

于 $x_2$ 方向的可变性.因此,大的 $x_1$ 坐标(用绝对值)不像大的 $x_2$ 坐标那样意外.当计算到原点的“距离”时,为 $x_2$ 坐标加权比为相同数值的 $x_1$ 坐标加权大一些,看起来是合理的.

接下来采用的方法是用样本标准差去除每个坐标.因此,被标准差除之后,得到“标准化”坐标 $x_1^* = x_1/\sqrt{s_{11}}$ 和 $x_2^* = x_2/\sqrt{s_{22}}$.现在这些标准化的坐标相互是在一个等同的基础上.在考虑变异性方面的不同后,我们利用标准欧几里德公式确定距离.

因此,从原点 $O = (0,0)$ 到点 $P = (x_1,x_2)$ 的统计距离可以由它的标准化坐标 $x_1^* = x_1/\sqrt{s_{11}}$ 和 $x_2^* = x_2/\sqrt{s_{22}}$ 计算出来.

\phantomsection\label{eq:1-13}
\begin{equation}
\begin{aligned}
d(O,P) &= \sqrt{(x_1^*)^2 + (x_2^*)^2} \\
&= \sqrt{\left(\frac{x_1}{\sqrt{s_{11}}}\right)^2 + \left(\frac{x_2}{\sqrt{s_{22}}}\right)^2} = \sqrt{\frac{x_1^2}{s_{11}} + \frac{x_2^2}{s_{22}}}
\end{aligned} \tag{1-13}
\end{equation}

比较\hyperref[eq:1-13]{式(1-13)}和\hyperref[eq:1-9]{式(1-9)},我们发现两式的不同是由于\hyperref[eq:1-13]{式(1-13)}中为 $x_1^2$ 和 $x_2^2$ 增加了权数 $k_1 = 1/s_{11}$ 和 $k_2 = 1/s_{22}$.注意,若样本方差是相同的, $k_1 = k_2$,则 $x_1^2$ 和 $x_2^2$ 将得到相同的权数.当权数相同时,忽略公因子,使用普通的欧几里德距离公式是很合适的.换言之,如果 $x_1$ 方向的变异性与 $x_2$ 方向的相同,并且 $x_1$ 值的变化不依赖于 $x_2$ 的值,欧几里德距离是适用的.

使用\hyperref[eq:1-13]{式(1-13)},我们看到所有到原点的距离的平方为常数 $c^2$ 的点,设坐标为 $(x_1,x_2)$,都满足

\phantomsection\label{eq:1-14}
\begin{equation}
\frac{x_1^2}{s_{11}} + \frac{x_2^2}{s_{22}} = c^2 \tag{1-14}
\end{equation}

\hyperref[eq:1-14]{式(1-14)}是一个椭圆方程,该椭圆以原点为中心,长轴和短轴分别与坐标轴重合.也就是说,\hyperref[eq:1-12]{式(1-12)}中统计距离的含义是,到原点距离为某常数的所有点的轨迹构成一椭圆.\hyperref[fig:1.21]{图 1.21} 表示了这种一般情况.

\begin{mvfigure}
\mvinc{pg-036}{0.119,0.399,0.454,0.529}
\phantomsection\label{fig:1.21}\mvcap{图 1.21\quad 常数统计距离的椭圆}
\end{mvfigure}

$d^2(O,P) = x_1^2/s_{11} + x_2^2/s_{22} = c^2$

\mvbold{例 1.14}\quad （计算统计距离）

两个变量的一组测量值 $(x_1,x_2)$ 得出 $\bar{x}_1 = \bar{x}_2 = 0, s_{11} = 4$ 和 $s_{22} = 1$.假设测量值 $x_1$ 与测量值 $x_2$ 无关,即一对测量值中任一个量的变化都不依赖于另一个.由于样本方差不相等,我们用

\[
d^2(O,P) = \frac{x_1^2}{4} + \frac{x_2^2}{1}
\]

度量任意点 $P = (x_1,x_2)$ 到原点 $O = (0,0)$ 的距离平方.所有到原点固定距离为 1 的点 $(x_1,x_2)$ 满足方程

\[
\frac{x_1^2}{4} + \frac{x_2^2}{1} = 1
\]

距离原点 1 个单位的一些点的坐标如下表所示：

\begin{mvtable}
\begin{tabular}{ll|l|l}
\multicolumn{2}{c|}{坐标: $(x_1,x_2)$} & \multicolumn{2}{c}{距离: $\dfrac{x_1^2}{4} + \dfrac{x_2^2}{1} = 1$} \\
\hline
$(0,1)$ & & $\dfrac{0^2}{4} + \dfrac{1^2}{1} = 1$ & \\
$(0,-1)$ & & $\dfrac{0^2}{4} + \dfrac{(-1)^2}{1} = 1$ & \\
$(2,0)$ & & $\dfrac{2^2}{4} + \dfrac{0^2}{1} = 1$ & \\
$(1,\sqrt{3}/2)$ & & $\dfrac{1^2}{4} + \dfrac{(\sqrt{3}/2)^2}{1} = 1$ & \\
\end{tabular}
\end{mvtable}

方程 $x_1^2/4 + x_2^2/1 = 1$ 的图形是一个中心在 $(0,0)$ 的椭圆,它的长轴沿着 $x_1$ 坐标轴,短轴则
